% ---------------------------------------------------------------------
% ANNEXE : pièces justificatives.
% Regroupe le matériel dont le corps du mémoire a besoin pour être
% vérifiable mais pas pour être compris : réconciliation des chiffres
% entre chapitres, contrôles de la conversion de sévérité, inventaire
% des paramètres, sensibilités hors tornado.
% ---------------------------------------------------------------------
\chapter{Pièces justificatives}
\label{ann:pieces}

Cette annexe rassemble ce qui rend le mémoire \emph{vérifiable} sans être nécessaire à sa
lecture : le dispositif qui contrôle chaque nombre publié, la réconciliation des chiffres entre
chapitres, deux contrôles de la brique de sévérité, l'inventaire complet des paramètres et les
sensibilités qui ne figurent pas au tornado. Aucun chiffre du corps n'y est établi pour la
première fois, et aucune conclusion n'en dépend ; l'objet est qu'aucun chiffre du mémoire ne soit
sans origine déclarée.

\section{Le dispositif de vérification des nombres publiés}
\label{sec:ann-harnais}
% HARNAIS-HORS-SECTION: tous les nombres de cette section decrivent l'INSTRUMENT et non le
% modele. Les faire verifier par lui serait circulaire. La declaration est DANS la section et
% non avant son titre, et la section ne cite aucun script : le harnais n'honore une declaration
% que si la section est par ailleurs sans source.

Un mémoire de cette taille publie plusieurs milliers de nombres, et le risque n'est pas qu'un
calcul soit faux : c'est qu'un nombre juste au moment où il a été écrit cesse de l'être quand une
valeur amont bouge, sans que personne ne s'en aperçoive. Ce mémoire s'est donc doté d'un harnais
de vérification, dont la portée se décrit ici avec ce qu'il ne garantit pas.

\paragraph{Le principe.} Pour chaque section, le harnais relève les scripts qu'elle cite, puis
cherche chacun des nombres qu'elle publie dans les \emph{sorties versionnées} de ces seuls
scripts. Un nombre est confirmé s'il y figure à la tolérance d'arrondi près,
\[
  \max\bigl(0{,}6\,\% \;;\; \text{demi-unité du dernier chiffre écrit}\bigr),
\]
soit $0{,}05$ pour un nombre écrit \og $2{,}1$ \fg{} et $0{,}5$ pour \og $122$ \fg{}. Une tolérance à plancher absolu,
telle qu'elle existait d'abord, faisait confirmer $0{,}9$ par un $0{,}99$ sans rapport.

\paragraph{Il rend trois lignes, et pas une.} Vérifiés, exemptés \emph{par motif}, non confirmés.
Une exemption n'est jamais silencieuse, faute de quoi le dénominateur ne serait plus \og les
nombres publiés \fg{} mais \og les nombres publiés moins une liste \fg{}, sans dire laquelle. Les
exemptions subsistantes sont les niveaux de confiance, les entiers d'énumération et les
millésimes. Une exemption se juge sur le contexte, jamais sur la valeur : c'est en indexant sur la
valeur qu'un résultat central avait été retiré du contrôle.

\paragraph{La couverture passe avant le taux.} Une section qui ne cite aucun script est \emph{hors
contrôle} : ses nombres ne sont ni confirmés ni infirmés, et cela ne se voit pas dans un taux. Un
taux se règle en retirant une citation, une couverture non. Le contrôle de fin de travail est donc,
dans cet ordre, couverture égale à $100\,\%$ puis confirmation au moins égale à $97\,\%$. Le fait
qui a imposé cet ordre est empirique : le périmètre a d'abord été de $74{,}7\,\%$ des nombres
publiés, et les valeurs périmées retrouvées lors de son extension venaient toutes de la zone non
couverte ou d'une exemption silencieuse, jamais d'un non confirmé. Trois états sont donc
distingués, sous contrôle, déclaré hors script avec son motif, et hors contrôle : seul le dernier
doit valoir zéro.

\paragraph{Ce que le harnais ne fait pas.} Il ne vérifie pas qu'un nombre est au bon endroit, ni
qu'il veut dire ce que la phrase prétend : un nombre confirmé peut être mal commenté. Il ne détecte
pas non plus les contradictions entre scripts, idée instruite puis écartée, un détecteur de
quasi-collision sur les valeurs confirmant par hasard dans un pool de plusieurs milliers de
nombres. Ce qui la remplace est une convention de rédaction : toute grandeur simulée se publie avec
son bruit. Le résidu se répartit en quatre classes dont une seule est fautive, la grandeur citée
mais jamais imprimée, d'où la règle appliquée partout : tout rapport cité doit être imprimé par un
script.

\paragraph{Ce que le dispositif a trouvé.} Trois erreurs réelles, qu'aucune relecture n'avait
vues : une borne d'intervalle de l'indice de queue arrondie dans le mauvais sens sur deux
chapitres ; un estimateur de Hill publié à une valeur qu'aucune population du pipeline ne
reproduit, corrigée au chapitre~\ref{chap:socle} ; et un seuil d'inefficacité publié à deux valeurs
distinctes à vingt-cinq lignes d'écart. Chacune est une valeur juste au moment où elle a été
écrite, devenue fausse quand la valeur amont a bougé, et seul un contrôle systématique pouvait les
révéler.

\begin{cle}
\textbf{État du contrôle}, tous chapitres, à la dernière exécution datée : 2\,208 nombres publiés
et confirmés, pour une couverture de $100\,\%$, aucun nombre n'étant hors contrôle sans
déclaration motivée. Le dispositif s'appuie sur 103 sorties de scripts versionnées et tourne aussi
sur les supports de présentation, où il a déjà rattrapé une valeur fausse.

\textbf{Pourquoi cette section échappe au contrôle qu'elle décrit.} La faire vérifier par le
harnais serait circulaire : il confirmerait le nombre qui compte ce qu'il confirme. Elle est donc
déclarée hors script avec ce motif, et ses valeurs sont celles d'une exécution datée, rejouable en
une commande.
\end{cle}

\section{Réconciliation du socle et de l'état conforme}
\label{sec:ann-reconciliation}

Le chapitre~\ref{chap:resultats} parle d'un SCR \emph{conforme} d'environ $5\,900$~M\euro{},
quand le chapitre~\ref{chap:partielle} parle d'un \emph{socle} de $5\,275$~M\euro{}. Ce ne
sont pas deux mesures de la même chose, mais deux points d'un seul modèle. Le socle est pris à
contagion nulle ($W=0$), fréquence de référence, sans canal de détection. L'état
conforme, lui, conserve une contagion (le gain $g=0{,}45$ n'est pas nul) et module la
détection. Le pont, calculé sur le moteur commun à graine fixée, l'établit terme à terme :
\[
  \underbrace{5\,275}_{\text{socle }W=0}\;\xrightarrow{+\text{ contagion }g=0{,}45}\;6\,504
  \;\xrightarrow{+\text{ détection }\times 0{,}85}\;5\,936\;\approx\;\text{conforme}.
\]
Le canal fréquence est neutre ici, le scénario conforme S0 ayant la fréquence de référence.
L'écart résiduel entre le moteur exact du chapitre~\ref{chap:partielle} et le Monte-Carlo du
chapitre~\ref{chap:resultats} est de $210$~M\euro{} ($4\,\%$) sur une configuration identique,
soit le bruit d'échantillonnage déjà déclaré, non une divergence de fond. L'ordre socle $<$
conforme $<$ non conforme est donc cohérent : le socle est plus bas \emph{parce qu'}il n'a
aucune contagion.

\begin{figure}[htbp]\centering
\figover{Z7_reconciliation.png}
\caption{Réconciliation des chiffres. (a)~Du socle ($W=0$, chapitre~\ref{chap:partielle}) à
l'état conforme (chapitre~\ref{chap:resultats}), un canal ajouté à la fois sur le moteur
commun. (b)~Sur une configuration identique, l'écart entre le moteur exact et le Monte-Carlo
est le seul bruit d'échantillonnage.}
% SOURCES-SCRIPTS: 36 67
\label{fig:reconciliation}
\end{figure}

\section{L'intervalle PRC tient quand on cesse de croire la conversion exacte}
\label{sec:ann-jacobs}

La sévérité PRC est dérivée du nombre d'enregistrements compromis par la relation log-log de
Jacobs (2014), dont la sortie de régression donne un écart-type résiduel de $0{,}523$ en
logarithme naturel ($R^2 = 0{,}51$) : à volume donné, le coût réel varie d'un facteur $2{,}4$
autour de la droite à $90\,\%$. Réinjecter ce résidu enregistrement par enregistrement dans le
bootstrap ne rouvre pourtant presque pas l'intervalle (facteur d'IC $\approx 1{,}4$ avec comme
sans résidu) : sur $15\,053$ enregistrements, le bruit se moyenne dans l'ajustement GPD, et le
plafond de sévérité borne la queue épaissie. L'effet est un déplacement de \emph{niveau} : la
queue dérivée s'épaissit ($\xi$ de $1{,}03$ à $1{,}07$, taux de matérialité de $0{,}150$ à
$0{,}157$) et la médiane du surcoût monte de $7\,\%$ (de $3\,600$ à $3\,861$~M\euro).
L'étroitesse de l'intervalle PRC relevée au chapitre~\ref{chap:resultats} n'est donc pas un
artefact de la conversion déterministe. L'incertitude qui ne se moyenne pas, elle, est celle
des coefficients $(a,b)$ de la régression, estimés sur $115$ observations : elle relève d'un
axe de sensibilité, traité au \S\ref{sec:ann-hypotheses}, non d'un bruit par enregistrement.

\begin{figure}[htbp]\centering
\figover{S16_prc_jacobs_residu.png}
\caption{Le résidu de la conversion Jacobs ($\mathrm{RSE} = 0{,}523$) : la distribution du
surcoût PRC se déplace plus qu'elle ne s'élargit.}
% SOURCES-SCRIPTS: 21 67
\label{fig:jacobsresidu}
\end{figure}

\section{Tableau des paramètres non calibrés et de leur sensibilité}
\label{sec:ann-parametres}

Trois natures sont distinguées : calibré sur données (avec son incertitude),
\emph{posé} par jugement d'expert ou choix normatif (non identifiable, cf.\
chapitre~\ref{ch:identifiabilite}), et semi-ancré sur un proxy externe non
spécifique au périmètre.

\vspace{0.25cm}
% PASSEE DE tabular A longtable LE 17 AOUT 2026. Cette table tenait sur une page tant que rien
% ne la precedait dans le chapitre ; l'ajout de la section sur le dispositif de verification l'a
% decalee et elle a debordé de 102 pt, avec l'annotation hyperref hors page qui accompagne
% toujours ce debordement. Un tabular NE PEUT PAS se couper entre deux pages : le piege est
% documente et il vient de mordre. longtable le peut, et l'en-tete se repete.
{\footnotesize
\renewcommand{\arraystretch}{1.28}
\begin{longtable}{@{}>{\raggedright\arraybackslash}p{2.7cm}
                    >{\raggedright\arraybackslash}p{2.6cm}
                    >{\raggedright\arraybackslash}p{1.7cm}
                    >{\raggedright\arraybackslash}p{4.0cm}
                    >{\raggedright\arraybackslash}p{3.6cm}@{}}
\toprule
\textbf{Paramètre} & \textbf{Valeur} & \textbf{Nature} & \textbf{Justification} & \textbf{Sensibilité du SCR} \\
\midrule
\endfirsthead
\toprule
\textbf{Paramètre} & \textbf{Valeur} & \textbf{Nature} & \textbf{Justification} & \textbf{Sensibilité du SCR} \\
\midrule
\endhead
\bottomrule
\endlastfoot
Queue de sévérité $\xi$ & $0{,}5954$ (OpRisk, chiffre de tête) ; PRC $1{,}03$ ; $0{,}90$ posé en pire cas de modèle & calibré, incertain & GPD sur pertes réelles ; hétérogène par catégorie ($0{,}87$--$1{,}37$) & \textbf{forte} : de $0{,}60$ à $0{,}38$ quand le seuil monte au-dessus de celui retenu, dans l'IC90 $[0{,}30\,;0{,}83]$ ; pilote la queue \\
Échelle GPD $\sigma$, seuil $u$, $p_u$ & $57{,}97$ M\euro{} ; $20{,}03$ M\euro{} ; $0{,}15$ & calibré & POT au percentile $84{,}4$ ; $p_u$ gelé et en écart de $+3{,}6\,\%$ avec ce seuil, effet chiffré au \S\ref{sec:pu-coherence} & modérée (via la queue) \\
Fréquence $\lambda$ & ${\sim}21{,}6$/an (OpRisk) & calibré & incidents TIC matériels / 27 ans & linéaire sur le niveau \\
Multiplicateur de fréquence par état & $\times 1{,}525$ partiel ; $\times 2{,}487$ non conforme ($\lambda$ de $21{,}56$ à $53{,}62$) & semi-ancré & parts Hackmageddon et fourchettes posées sur des effets publiés (chapitre~\ref{chap:multietats}) ; aux bornes, $\lambda$ de $41{,}7$ à $65{,}6$ & \textbf{forte} : canal dominant du chiffre de tête \\
Sur-dispersion $\phi$ & $9{,}20$ & \textbf{posé} (prudent) & $2{,}24$ mesuré sur la série annuelle détendancée, maintenu entre scénarios & faible sur le niveau \\
Choc commun fréquence $a$ & $0{,}60$ & \textbf{posé} (prudent) & dépasse l'IC empirique haut ($0{,}22$) : conservateur & enveloppe le modèle à paquets ($\times 2{,}5$) \\
Structure $W = g\,\texttt{TRANS}/\max_k s_k$ & classeur qualitatif & \textbf{posé} / normatif & jugement d'expert consigné ; direction non identifiable & l'asymétrie départage les deux piliers de tête ; l'essentiel du surcoût de contagion vient de la partie symétrique \\
Gain de propagation $g$ & $0{,}90$ base ; par état $0{,}45/0{,}68/0{,}90$ & \textbf{posé} & amplitude de $W$, bornée par Leontief $\rho(W)\le g<1$ & \textbf{levier} : $\times 1{,}6$ sur $[0{,}05\,;0{,}90]$ \\
Score de conformité $q$ / états & continu $[0,1]$ (ou NC/PC/C) & \textbf{posé} & décrit le profil, subsume 3 ou 5 niveaux & \textbf{levier} : $-53\,\%$ de NC à C sur le seul score ; facteur $3{,}34$ sur les quatre canaux \\
Multiplicateur de matérialité par état & $0{,}85/1{,}00/1{,}20$ & \textbf{posé} & module $p_u$ selon la conformité & modérée (canal détection) \\
Charge systémique $\gamma$ & $0{,}68$ & semi-ancré & proxy concentration cloud (AWS+Azure+GCP, Synergy 2025) & modérée \\
Seuil de crise $\theta_{\text{crise}}$ & $-2{,}5$ & \textbf{posé} & déclenche le régime de co-mouvement extrême & agit sur la seule queue crise \\
Probas d'état $P(\text{NC/PC/C})$ & $0{,}35/0{,}35/0{,}30$ & semi-ancré & enquêtes conformité DORA ${\sim}50\,\%$ non conformes (Deloitte, ESA) & déplace le mélange, pas les états \\
Copule de dépendance $\theta$ & Gumbel $1{,}8$ & \textbf{posé} & queue supérieure entre briques & \textbf{négligeable} : $\Delta_{\text{DORA}}$ $<1{,}4\,\%$ sur $[1,4]$ \\
Excès de source $u_{ij}$ & $\operatorname{logit}p_{ij}=\operatorname{logit}p_j+u_{ij}$ & \textbf{non identifié} & la direction elle-même ; normalisation $\sum_i u_{ij}=0$ (repérage, non information) & mesurée par l'échelle de repli (ch.~\ref{chap:partielle}) : rang~$1$ capte $74\,\%$ et resserre les bornes de $25\,\%$ ; à $u=0$ P1 et P4 deviennent indiscernables \\
Élasticité sévérité/taille $b$ & $0{,}087$ $[0{,}026\,;0{,}148]$ & calibré, tronqué & EMV lognormale tronquée au seuil de collecte (ch.~\ref{chap:donnees}) & faible : $\times 0{,}86$ sur la sévérité transposée \\
Plafond de sévérité & $40$ M\euro{} & \textbf{posé} & capacité de réassurance ($\xi\ge 1$) & borne la queue extrême \\
\end{longtable}
}

\vspace{4pt}
% SOURCES-SCRIPTS: 08b 22 25 34 47 60 67 99 43

\textbf{Ce que la table rend visible.} Les seuls paramètres à forte sensibilité sont
soit calibrés sur données (la queue $\xi$, dont on assume et chiffre l'incertitude), soit des
\emph{leviers voulus} ($g$, $q$, le multiplicateur de fréquence) qui traduisent précisément l'objet du mémoire, l'effet de
la conformité et de la contagion sur le capital. Les paramètres purement posés qui
ne sont pas des leviers ($\theta$ de copule, $\gamma$, $\theta_{\text{crise}}$, probas d'état,
matérialité) ont une influence faible à modérée et documentée. Aucune conclusion ne
repose sur un paramètre à la fois non calibré \emph{et} déterminant hors de son rôle attendu.

\section{Les hypothèses hors tornado et leur axe de stress}
\label{sec:ann-hypotheses}

Trois entrées ne figuraient pas au tornado du chapitre~\ref{chap:resultats} : les coefficients
de la conversion Jacobs, la charge systémique $\gamma$ et l'ancrage des probabilités d'état.
Les coefficients $(a,b)$, estimés sur $115$ observations, sont stressés à $\pm 1{,}645$
écart-type en \emph{pivotant la droite au centroïde} de la régression, déduit des erreurs-types
publiées ($\bar{x} = 10{,}04$, soit environ $23\,000$ enregistrements) ; stresser $a$ et $b$
indépendamment ignorerait leur corrélation négative. L'élasticité $b$ commande l'épaisseur de
la queue dérivée ($\xi$ de $0{,}85$ à $1{,}21$) mais le surcoût PRC reste contenu, de $3\,125$
à $3\,977$~M\euro{} autour du central ($3\,577$), le plafond de sévérité bornant l'effet :
c'est l'axe d'hypothèse dominant du périmètre PRC, et le verdict y survit. $\gamma$ ne touche,
par construction, ni le SCR par état ni le SCR espéré en régime normal (invariant à
$10\,644$~M\euro) : il ne commande que la sévérité de la bascule en crise
($\mathrm{P}(\text{NC} \mid \text{crise})$ de $84$ à $100\,\%$, SCR espéré en crise de
$14\,148$ à $15\,071$). L'ancrage des probabilités d'état, enfin, déplace le capital espéré
actuel ($9\,951$ à $11\,485$~M\euro{} selon les enquêtes retenues) mais jamais le SCR par état
ni le surcoût NC contre C, qui lui sont conditionnels. Les trois axes confirment la ligne de
partage du mémoire : les hypothèses bougent des niveaux, jamais la direction ni le classement.
% SOURCES-SCRIPTS: 22 67

\begin{figure}[htbp]\centering
\figover{S17_sensibilites_hypotheses.png}
\caption{Sensibilités des hypothèses restantes : coefficients Jacobs (pivotés au centroïde
déduit des erreurs-types publiées), charge systémique $\gamma$, ancrage des probabilités
d'état.}
% SOURCES-SCRIPTS: 22
\label{fig:hypotheses}
\end{figure}

% =====================================================================
\section{Portée et angles morts de chaque proxy}
\label{sec:ann-proxys}

% TABLE DEMANDEE COMME INDISPENSABLE AU POINT D'ETAPE DU 13 AOUT : « quel que soit le choix de
% modelisation, il est indispensable d'expliciter ce que les proxys retenus capturent et ce qu'ils
% ne capturent pas, et d'assumer les choix effectues ».
%
% ELLE EST VOLONTAIREMENT PRESQUE SANS CHIFFRES, et c'est un choix de conception. Sa valeur est
% dans les ENONCES : chaque nombre vit deja sous le script qui le produit, et le repeter ici
% multiplierait les occurrences sans rien ajouter, tout en gonflant le risque de confirmation
% fortuite que le harnais documente. Les renvois font le lien.

\noindent Un modèle qui manque de données procède par \emph{proxys}. Le danger n'est pas d'en
utiliser, il est de laisser croire qu'ils mesurent ce qu'ils remplacent. Cette table les recense
et dit, pour chacun, ce qu'il porte et ce qu'il ne porte pas. Elle se lit comme la contrepartie de
la table des paramètres du \S\ref{sec:ann-parametres} : celle-là dit ce qui est calibré ou posé,
celle-ci dit sur \emph{quoi}.

\vspace{0.2cm}
{\footnotesize
\renewcommand{\arraystretch}{1.3}
% COLONNES AU FER A GAUCHE ET NON JUSTIFIEES. Sur 3 cm, la justification ecarte les mots jusqu'a
% produire des blancs de la largeur d'un caractere au milieu d'un terme (« Classeur   qualitatif »).
% Defaut vu en regardant la page rendue, invisible a la compilation.
\begin{longtable}{@{}>{\raggedright\arraybackslash}p{3.1cm}
                    >{\raggedright\arraybackslash}p{5.4cm}
                    >{\raggedright\arraybackslash}p{5.6cm}@{}}
\toprule
\textbf{Proxy} & \textbf{Ce qu'il capture} & \textbf{Ce qu'il ne capture pas} \\
\midrule
\endfirsthead
\toprule
\textbf{Proxy} & \textbf{Ce qu'il capture} & \textbf{Ce qu'il ne capture pas} \\
\midrule
\endhead
\bottomrule
\endlastfoot
Classeur qualitatif ($\mathrm{TRANS}$, $\mathrm{ROOT}$) &
L'\emph{ordre} d'influence entre piliers, consigné et traçable ; il fournit la structure que
la donnée ne donne pas &
Aucune magnitude vérifiable, et aucune direction identifiable : c'est un jugement, non une
mesure. Traité par identification partielle, ch.~\ref{chap:partielle} \\

Chaîne de saut inter-piliers &
La direction des transferts, et le courant qui la mesure ($\S\ref{sec:ann-martingalisation}$) &
Ni l'extinction, ni l'épuisement des piliers déjà touchés : c'est une \emph{projection}
markovienne d'un processus auto-évitant à $160$ états \\

Corpus de post-mortems &
L'\emph{ordre} de défaillance dans des incidents documentés, testé contre une loi nulle exacte &
La magnitude : le corpus établit un ordre, pas une intensité. Et il est sujet au biais de
narration, mesuré au \S\ref{sec:biais-narration} \\

OpRisk Global, cyber $\times$ finance &
La queue des pertes opérationnelles matérielles au-delà du seuil de collecte &
L'attritionnel, absent par troncature ; et l'exposition informatique propre à une entité \\

Chronologie de brèches (PRC) &
Les comptages et les \emph{dates}, donc la structure temporelle et la concentration &
Le coût, qui n'y figure pas : il est \emph{dérivé} par la conversion ci-dessous \\

Conversion log-log de Jacobs &
La relation moyenne entre enregistrements compromis et coût &
Le coût d'un incident donné : le résidu vaut un facteur $2{,}4$ à $90\,\%$. Il se moyenne dans
l'ajustement, ce qui est vérifié au \S\ref{sec:ann-jacobs}, mais ne disparaît pas \\

Concentration journalière des tiers &
La concentration \emph{en coupe} sur des prestataires partagés, une cause frappant de nombreuses
entités le même jour &
L'accumulation intra-annuelle \emph{d'une} entité : les victimes d'un jour sont des entités
distinctes, non les sinistres d'une seule \\

Concentration du marché de l'informatique en nuage &
Un ordre de grandeur pour la charge systémique, ancré sur une source publique &
La dépendance \emph{propre} de l'entité à ces fournisseurs, que seul le registre
d'externalisation donnerait \\

Enquêtes de conformité DORA &
Une part de marché de non-conformité, pour ancrer les probabilités d'état &
La répartition par pilier et par entité : l'ancrage est global, l'usage est par pilier \\

Facteur systémique commun &
Le co-mouvement des états de conformité par l'environnement, avec une charge forte &
La \emph{structure} de ce co-mouvement, échangeable ici alors que le recouvrement des contrôles
est spécifique. Effet mesuré au \S\ref{sec:dependance-etats} \\

Élasticité sévérité/taille &
La relation moyenne entre taille et sévérité sur le panel, estimée par vraisemblance tronquée &
Un plafond à l'échelle d'entité : l'élasticité contient presque zéro, d'où la borne inférieure
de validité et la requalification du chiffre d'entité en borne supérieure \\

Bilans \textsc{sfcr} &
La taille réelle de quatre entités, sur données réglementaires publiées &
Leur exposition informatique, et leur état de conformité, qui est \emph{supposé} et non observé \\

Hackmageddon &
La \emph{structure} du risque par vecteur d'attaque, c'est-à-dire des parts &
Le niveau, jamais. Et la reproductibilité : le jeu a été consulté puis non conservé, ce que le
mémoire déclare comme citation externe non recalculable \\

Marge de risque réglementaire &
Le coût de détention du capital tel que le régime prudentiel le définit &
Le coût du capital propre à l'entité. Et le taux a changé, ce qui est chiffré plutôt que
corrigé au \S\ref{sec:roi} \\
\end{longtable}
}

\vspace{0.2cm}
\noindent Trois lectures de cette table méritent d'être dégagées, parce qu'elles ne sont visibles
qu'une fois les proxys mis côte à côte.

\textbf{La colonne de droite n'est pas une liste de regrets.} Chaque case y est soit traitée
ailleurs, soit déclarée comme limite, et les renvois le montrent. Ce qui serait fautif ne serait
pas d'avoir des proxys imparfaits, ce serait de ne pas savoir lesquels.

\textbf{Deux proxys se relaient plutôt qu'ils ne se complètent}, et il faut le voir : la
chronologie de brèches porte les dates mais pas les coûts, la conversion porte les coûts mais pas
les dates. Le coût d'un incident \emph{daté} n'est donc jamais observé, il est toujours
reconstruit. C'est la faiblesse la plus structurelle de la donnée disponible, et aucun des deux
proxys ne la révèle seul.

\textbf{Et le même mot, \og concentration \fg{}, désigne deux objets dans cette table.} La concentration
journalière est en coupe, entre entités ; l'accumulation que le modèle fait porter au pilier des
tiers est au sein d'un incident, entre piliers. Le premier calibre le second par analogie, non par
mesure, et c'est une hypothèse à assumer plutôt qu'un ancrage.

% SOURCES-SCRIPTS: 34 40 57 59 63 69

% =====================================================================
\section{La CTE au niveau agrégé, et à quel niveau elle égale la VaR réglementaire}
\label{sec:ann-cte}

% SOURCES-SCRIPTS: 73

Le chapitre~\ref{chap:resultats} établit que la mesure de queue sert à \emph{mesurer} et non à
\emph{couvrir}. Cette pièce en donne les nombres, qui manquaient : la CTE n'était calculée que sur
le quantile d'\emph{un} sinistre, jamais sur la charge annuelle, et son argument d'existence était
invoqué plutôt que montré.

\paragraph{La première réponse est réglementaire, et il faut la donner en premier.} Le régime
applicable définit le capital de solvabilité requis comme une valeur en risque à
$99{,}5\,\%$ sur un horizon d'un an. Un travail qui reporterait une mesure de queue à la place
répondrait à une autre question que celle posée, indépendamment de ses mérites : le Test suisse
de solvabilité retient une espérance conditionnelle de queue à $99\,\%$, et le débat entre les
deux régimes est ancien, documenté et tranché différemment selon les juridictions. Ce mémoire se
place sous le premier régime, et c'est la raison principale pour laquelle il ne substitue pas
l'une à l'autre. Les arguments qui suivent viennent \emph{renforcer} ce choix, ils ne le
fondent pas : les présenter dans l'ordre inverse laisserait croire qu'on ignore le cadre.

\subsection*{La CTE agrégée, et le niveau auquel elle égale la VaR réglementaire}

\begin{center}\footnotesize
\renewcommand{\arraystretch}{1.25}
\begin{tabular}{@{}l r r r@{}}
\toprule
\textbf{Niveau $\beta$} & \textbf{VaR} & \textbf{CTE} & \textbf{CTE\,/\,VaR} \\
\midrule
$95\,\%$ & $3\,055$ & $5\,734$ & $1{,}88$ \\
$99\,\%$ & $6\,078$ & $12\,607$ & $2{,}07$ \\
$99{,}5\,\%$ & $\mathbf{8\,374}$ & $18\,183$ & $2{,}17$ \\
\bottomrule
\end{tabular}
\end{center}

\vspace{4pt}
\noindent En M\euro, périmètre OpRisk, calibration figée. Le nombre décisif n'est aucun de ceux-là
mais celui qui permet de \emph{comparer} les deux conventions :
\begin{equation}
  \mathrm{CTE}_\beta(L) \;=\; \VaR_{99{,}5\%}(L) \quad\text{pour}\quad \beta = 97{,}73\,\%.
  \label{eq:cte-equivalent}
\end{equation}
Une CTE à $95\,\%$ vaut donc $5\,734$~M\euro{} contre $8\,374$ : rapportée comme capital, elle
serait moins prudente que l'exigence de Solvabilité~II, dans un rapport de $1{,}46$. Sans
\eqref{eq:cte-equivalent}, \og CTE à $95\,\%$\fg{} et \og VaR à $99{,}5\,\%$\fg{} ne se comparent pas et l'on
ne sait pas laquelle est la plus exigeante ; avec elle, l'arbitrage devient arithmétique. C'est la
raison pour laquelle la CTE est publiée en diagnostic : substituer une mesure à l'autre
déplacerait le niveau d'exigence de moitié sans que ce soit une décision.

Le rapport $\mathrm{CTE}/\VaR$ agrégé vaut $2{,}17$ à $99{,}5\,\%$, sous l'asymptote
$1/(1-\xi) = 2{,}47$ de la sévérité : l'agrégat s'en approche par le principe de la perte unique
dominante mais reste en deçà à niveau fini. C'est le \emph{rapport} qui informe, non les deux
montants.
% SOURCES-SCRIPTS: 73

\subsection*{L'existence, montrée plutôt qu'invoquée, et sur la bonne quantité}

Le protocole est simple : même moteur, \emph{sans} plafond, à deux indices de queue, et l'on
regarde la dispersion de la CTE entre six graines quand le nombre d'années croît.

\begin{center}\footnotesize
\renewcommand{\arraystretch}{1.25}
\begin{tabular}{@{}r r r r r@{}}
\toprule
\textbf{Années} & \textbf{CTE ($\xi=0{,}5954$)} & \textbf{Dispersion} & \textbf{CTE ($\xi=1{,}0328$)} & \textbf{Dispersion} \\
\midrule
$50\,000$ & $18\,084$ & $7{,}4\,\%$ & $998\,016$ & $51{,}4\,\%$ \\
$200\,000$ & $18\,263$ & $4{,}5\,\%$ & $1\,203\,706$ & $44{,}4\,\%$ \\
$800\,000$ & $18\,023$ & $4{,}2\,\%$ & $1\,182\,736$ & $38{,}6\,\%$ \\
\bottomrule
\end{tabular}
\end{center}

\vspace{4pt}
\noindent \textbf{Ce qui tranche est le niveau de la dispersion, non sa vitesse.} À l'indice
publié, la CTE se cite à deux chiffres ; au-delà de $1$, elle reste entre $39$ et $51\,\%$ de sa
propre valeur, de sorte que deux échantillons de huit cent mille années en donnent des valeurs
qui diffèrent de deux cinquièmes. Elle ne se cite alors \emph{pas}, même à un chiffre
significatif, et le niveau atteint est sans interprétation : il mesure la taille du plus grand
tirage de l'échantillon, non un risque.

\begin{attention}
\textbf{La vitesse, elle, ne conclut pas, et c'est un résultat du modèle plutôt qu'une limite du
test.} Les dispersions sont multipliées par $0{,}57$ et $0{,}75$ sur seize fois plus d'années, là
où un estimateur classique donnerait $0{,}25$. \emph{Aucune} des deux séries n'atteint cette
référence, et c'est attendu : à $\xi = 0{,}5954$ la variance de la sévérité est déjà
infinie même si son espérance existe, donc aucun estimateur de ce modèle ne converge au rythme
classique. Cela vaut aussi pour la VaR, et c'est une raison de plus de publier les grandeurs
simulées avec leur bruit.
% SOURCES-SCRIPTS: 73
\end{attention}

\subsection*{Sur le périmètre à queue plus lourde, le plafond vide la mesure de son contenu}

Le périmètre PRC porte un indice de $1{,}0328$ et un plafond de sévérité de $40$~M\euro, qui rend
l'espérance finie artificiellement. En faisant varier ce plafond, on constate qu'il multiplie la
CTE et la VaR du \emph{même} facteur, $1{,}67$ contre $1{,}66$ pour un plafond quadruplé : il agit
comme une échelle sur toute la loi et ne distingue pas la CTE. L'information est donc dans le
rapport, et il vaut $1{,}05$ sous PRC contre $2{,}17$ sur la calibration publiée. Le plafond a
retiré la queue que la CTE est censée mesurer.
% SOURCES-SCRIPTS: 73

\begin{cle}
\textbf{L'alternative qui ferme la question du niveau de couverture.} Sur un périmètre d'indice
supérieur à $1$, \emph{ou bien} l'on ne plafonne pas et la CTE n'est pas reproductible, \emph{ou
bien} l'on plafonne et son rapport à la VaR tombe à $1{,}05$, de sorte qu'elle n'apporte plus rien
au-delà du quantile. Dans les deux cas elle ne peut pas porter le capital. Là où l'indice est
inférieur à $1$, elle est parfaitement légitime, et c'est l'usage qu'en fait ce mémoire : un
diagnostic de forme de queue, et un noyau d'allocation qui répartit un niveau déjà fixé.
\end{cle}

% =====================================================================
\section{Le niveau de l'intervalle reporté : de 90 à 95 pour cent}
\label{sec:ann-intervalle95}
% SOURCES-SCRIPTS: 72

Le niveau de confiance de l'intervalle reporté est une convention de report : il se choisit. Le
niveau retenu reste $90\,\%$, et cette pièce en donne l'alternative chiffrée plutôt qu'une
substitution. Aucun bootstrap n'est rejoué, ce qui donnerait deux valeurs du même intervalle : le
passage se fait par la même approximation normale que celle dont le moteur se sert déjà, les deux
demi-largeurs étant mises à l'échelle séparément par $1{,}96/1{,}645 = 1{,}1915$, ce qui conserve
l'asymétrie. La calibration ne bouge pas, donc tous les capitaux publiés non plus.

\begin{center}\footnotesize
\renewcommand{\arraystretch}{1.25}
\begin{tabular}{@{}l l l r r@{}}
\toprule
\textbf{Grandeur} & \textbf{IC $90\,\%$} & \textbf{IC $95\,\%$} & \textbf{$\times$ à 90} & \textbf{$\times$ à 95} \\
\midrule
Indice de queue $\xi$ & $[0{,}3044\,;\,0{,}8313]$ & $[0{,}2487\,;\,0{,}8765]$ & $2{,}73$ & $3{,}52$ \\
Quantile de sévérité (M\euro) & $[411{,}5\,;\,1\,037{,}1]$ & $[363{,}4\,;\,1\,108{,}8]$ & $2{,}52$ & $3{,}05$ \\
Capital agrégé (M\euro) & $[3\,042\,;\,27\,152]$ & $[2\,727\,;\,34\,942]$ & $8{,}93$ & $12{,}81$ \\
\bottomrule
\end{tabular}
\end{center}

\vspace{4pt}
\noindent L'élargissement du capital est très asymétrique : sa borne basse ne descend que de
$315$~M\euro{} quand la haute monte de $7\,790$. C'est la convexité du capital en $\xi$, et elle
s'annonce plutôt qu'elle ne se subit.

\textbf{Le changement de niveau ne corrige pourtant rien, et c'est le point décisif.} La couverture
\emph{réelle}, mesurée à $n$ fini sur $2\,000$ échantillons de $91$ excès, vaut
$86{,}8 \pm 0{,}8\,\%$ pour un niveau annoncé de $90$ et $91{,}8 \pm 0{,}6\,\%$ pour un niveau
annoncé de $95$ : le manque est le même, $-3{,}2$ points. Relever le nominal déplace l'annonce sans
corriger l'estimateur, et la constance du déficit en désigne la cause, un écart-type asymptotique
trop petit à quatre-vingt-onze excès. Ce déficit se publie, parce qu'un intervalle annoncé à
$90\,\%$ qui en couvre $87$ engage celui qui l'annonce, mais il ne disqualifie pas l'intervalle :
à ce nombre d'excès, la convergence lente de la normalité asymptotique du maximum de vraisemblance
rend cet ordre de grandeur attendu.

\textbf{Le niveau reporté reste donc $90\,\%$.} Trois raisons convergent : le corps publie partout
les bornes de la configuration gelée, si bien qu'annoncer $95\,\%$ ici créerait une incohérence
pire que le défaut prétendument corrigé ; le passage à $95\,\%$ ne répare pas la sous-couverture ;
et le correctif qui la réparerait, un élargissement de la demi-largeur d'environ $10\,\%$, est une
recalibration que le gel interdit. La réserve du \S\ref{soc:sec:validation} reste due au nouveau
niveau comme à l'ancien.
% SOURCES-SCRIPTS: 63 72

% =====================================================================
\section{Martingalisation, tests conduits et tests impossibles}
\label{sec:ann-martingalisation}

% SOURCES-SCRIPTS: 40 64

Le mot \emph{martingalisation} a été employé dans les présentations orales de ce travail, et il
appelle une mise au point : un test de martingalité suppose que l'on identifie explicitement la
quantité testée, et cette identification n'avait pas été écrite. Cette pièce la fournit, en
séparant trois choses que le mot recouvrait.

\subsection*{La chaîne, et ses hypothèses mesurées plutôt que supposées}

L'objet n'est pas $W$ mais le noyau de saut inter-piliers de la cascade,
\[
  P_{jk} \;=\; \mathrm{TRANS}_{jk} \Big/ \sum_{k}\mathrm{TRANS}_{jk},
\]
soit la probabilité, sachant que la propagation quitte le pilier $j$, qu'elle aille vers $k$. Cette matrice est stochastique en
ligne, donc c'est un authentique noyau de Markov, ce que $W$ n'est pas. Les hypothèses qui
soutiennent tout ce qui suit sont mesurées et non postulées : la diagonale est nulle, les
\emph{vingt} entrées hors diagonale sont strictement positives, la chaîne est donc
\emph{irréductible} et apériodique, ce qui garantit l'existence et l'unicité de la loi
stationnaire $\pi = (0{,}131\,;\,0{,}263\,;\,0{,}200\,;\,0{,}192\,;\,0{,}214)$.
% SOURCES-SCRIPTS: 40

\subsection*{La martingale, nommée}

Pour une fonction test $f$ sur les cinq piliers, la décomposition de Dynkin
\begin{equation}
  M^f_r \;=\; f(X_r) - f(X_0) - \sum_{s=0}^{r-1}\bigl(Pf - f\bigr)(X_s)
  \label{eq:dynkin}
\end{equation}
est une martingale pour la filtration naturelle de la chaîne et sous la mesure
physique. Sa variation quadratique prévisible, en régime stationnaire, est la forme de Dirichlet
de l'\eqref{eq:dirichlet} du chapitre~\ref{ch:identifiabilite}. \emph{Martingaliser} désigne donc ici l'opération de
retirer le compensateur $\sum(Pf-f)$ pour ne garder que $M^f$, et l'énoncé du
chapitre~\ref{ch:identifiabilite} est que la variation quadratique ne dépend que de la partie
symétrique, le compensateur portant seul la direction.
% SOURCES-SCRIPTS: 40

\subsection*{Ce qui est effectivement testé, et ce n'est pas la martingale}

Aucune hypothèse n'est testée sur $M^f$, et il n'y en a pas à tester : \eqref{eq:dynkin} est une
\emph{identité algébrique}, vraie pour toute chaîne de Markov. Les tests statistiques de ce
volet portent sur un tout autre objet, la statistique d'asymétrie des transitions. Sur les
transitions annuelles d'OpRisk, confrontée à un placebo qui permute les années à l'intérieur de
chaque firme, elle ne détecte rien : l'asymétrie observée vaut $119$ contre $128 \pm 27$ sous le
nul, soit $z = -0{,}33$. Sur le corpus de post-mortems, la loi nulle est exacte, chaque arête y
changeant de sens indépendamment avec probabilité $\tfrac12$ (\S\ref{sec:biais-narration}). Aucun
de ces tests ne fait intervenir une martingale.
% SOURCES-SCRIPTS: 40 64

\subsection*{Pourquoi aucun test de martingalité n'est possible, et c'est la réponse de fond}

Un test de martingalité exige une trajectoire observée de la quantité candidate : on
vérifie que l'accroissement conditionnel est d'espérance nulle. Or les deux candidats de ce
mémoire, le facteur systémique de conformité et le facteur tiers, sont latents et jamais
observés : ils sont construits par le modèle, non mesurés. Il n'existe donc aucune série sur
laquelle un tel test porterait. Ce n'est pas un renoncement faute de temps, c'est une
impossibilité de principe, et elle se déclare.

Il faut ajouter que le cadre risque-neutre, auquel le mot \emph{martingalité} renvoie
habituellement, n'a pas lieu d'être ici : pas de prix, pas d'actif répliqué, pas de numéraire,
aucune hypothèse de non-arbitrage, aucun changement de mesure. Les tests usuels portent sur des
trajectoires actualisées au taux sans risque ; il n'y a dans ce mémoire ni actualisation ni
mesure équivalente.

\subsection*{Le déclenchement du pilier des tiers : la propriété utilisée n'est pas celle qu'on
nommait}

Au \S\ref{sec:p4-dynamique}, $Z$ est un mouvement brownien sans dérive, donc une
martingale. Mais ce n'est pas la propriété qui sert. Le principe de réflexion repose sur
la propriété de Markov forte et sur la symétrie des accroissements ; nommer la
martingale y était décoratif. Les hypothèses réellement requises sont : accroissements gaussiens
indépendants et stationnaires, absence de dérive, seuil $z^\star$ déterministe et fixé avant
l'observation, horizon déterministe. La réserve déjà publiée demeure : un brownien sans dérive
atteint le niveau presque sûrement mais avec une espérance de délai infinie, de sorte que la
médiane est la statistique reportable et non la moyenne.

\subsection*{L'hypothèse la plus lourde et les angles morts de la projection}

Elle n'est pas dans le vocabulaire mais dans l'objet. La chaîne $P$ est sans mémoire,
alors que la cascade effectivement simulée est auto-évitante : un pilier déjà tombé
éteint la propagation. Le processus réel n'est donc pas markovien sur les cinq piliers ; il l'est
sur le couple (pilier courant, ensemble déjà tombé), soit $5 \times 32 = 160$ états. L'analyse de
réversibilité porte par conséquent sur une projection markovienne du processus.

Il faut dire ce que cette projection capture et ce qu'elle ne capture pas, puisque c'est l'exigence
générale que ce travail s'est donnée. Elle capture la \emph{direction des transferts}, qui est
l'objet du chapitre, et le courant qui la mesure vaut $0{,}0443$ au maximum. Elle ne capture ni
l'\emph{extinction} de la cascade, ni l'\emph{épuisement} des piliers déjà touchés. Ces deux
mécanismes sont dans le moteur de simulation, ils ne sont pas dans la chaîne servant à l'argument
d'identification. L'argument par la réversibilité vaut donc pour les statistiques de transition
invariantes par renversement du temps ; la loi complète des ensembles de piliers touchés, que ces
mécanismes façonnent, dépend au contraire de la direction (proposition~\ref{prop:nonid}).
% SOURCES-SCRIPTS: 40

% =====================================================================
\section{La structure de la dépendance entre états et son absence d'effet}
\label{sec:ann-dependance-etats}

% SOURCES-SCRIPTS: 69 30

Cette pièce accompagne le \S\ref{sec:dependance-etats}, qui établit que la \emph{structure} de la
dépendance entre états de conformité ne déplace pas le capital espéré. La figure en donne les
deux moitiés de l'argument, et la seconde est celle qu'il faut regarder.

\begin{figure}[htbp]\centering
\figover{S25_structure_dependance_etats.png}
\caption{La dépendance des états : elle existe déjà, et sa structure n'agit pas.
\textbf{(a)}~le capital des $32$ configurations, moyenné par nombre de piliers non conformes,
est presque linéaire : les incréments vont de $2\,075$ à $2\,403$~M\euro, soit $15\,\%$ autour de
leur moyenne, et leur ordre n'est même pas monotone à cette résolution. (b)~l'écart au
capital espéré quand on ajoute un surcroît de corrélation aux deux paires concernées, tracé
\emph{à l'échelle de la largeur de la bande d'identification}. Les barres sont invisibles, et
c'est le résultat : $2$~M\euro{} au maximum contre $1\,839$ de largeur déclarée.}
% SOURCES-SCRIPTS: 69
\label{fig:dependance-etats}
\end{figure}

\noindent L'échelle du panneau~(b) fait partie du propos. Tracées sur leur propre amplitude, des
barres de $2$~M\euro{} paraîtraient considérables ; rapportées à l'incertitude que le mémoire
déclare déjà, elles disparaissent, et c'est l'échelle qui convient pour représenter une
invariance.

% =====================================================================
\section{Défaillances simultanées et coût de l'additivité}
\label{sec:ann-additivite}

% SOURCES-SCRIPTS: 74

Cette pièce accompagne le \S\ref{sec:additivite-couts}, qui distingue les trois énoncés
d'additivité du modèle et retient que celui portant sur les coûts d'un sinistre est le plus lourd
des trois. On en donne ici la mesure complète.

% La figure est declaree ICI, en tete de section, et non a sa place logique apres la seconde
% table : posee en fin de section elle partait seule sur une page de flottants aux trois quarts
% blanche, n'ayant plus aucun texte a partager.
\begin{figure}[!tb]\centering
\figover{S30_defaillances_simultanees.png}
\caption{Les défaillances simultanées et l'additivité de leurs coûts. (a)~la loi exacte
du nombre de piliers touchés, en échelle logarithmique pour que les cardinaux rares restent
visibles : la défaillance multiple devient majoritaire dès que la propagation s'ouvre.
\textbf{(b)}~le relâchement de l'additivité. L'écart parcourt $76\,\%$ de sa valeur publiée sans
jamais changer de signe, et les deux pentes diffèrent d'un facteur $2{,}42$ parce que l'exposant ne
mord que sur les sinistres multi-piliers, majoritaires au seul état non conforme.}
% SOURCES-SCRIPTS: 74
\label{fig:additivite-couts}
\end{figure}

\subsection*{La loi du nombre de piliers touchés, exacte plutôt que simulée}

Un sinistre tire un \emph{ensemble} de piliers dans la loi de progéniture du noyau auto-évitant.
Cette loi se calcule par énumération, sans aucun tirage, en pondérant les amorces par leur
propension. Elle est donc exacte, et c'est ce qui permet de la citer sans bruit associé.

\begin{center}\footnotesize
\renewcommand{\arraystretch}{1.25}
\begin{tabular}{@{}c c c@{}}
\toprule
\textbf{Piliers touchés par le sinistre} & \textbf{État conforme} ($g=0{,}45$) &
\textbf{Non conforme} ($g=0{,}90$) \\
\midrule
$1$ & $68{,}85\,\%$ & $37{,}69\,\%$ \\
$2$ & $25{,}10\,\%$ & $38{,}10\,\%$ \\
$3$ & $5{,}31\,\%$ & $18{,}29\,\%$ \\
$4$ & $0{,}70\,\%$ & $5{,}22\,\%$ \\
$5$ & $0{,}04\,\%$ & $0{,}69\,\%$ \\
\midrule
\emph{plus d'un pilier} & $\mathbf{31{,}15\,\%}$ & $\mathbf{62{,}31\,\%}$ \\
\emph{moyenne} & $1{,}380$ & $1{,}931$ \\
\bottomrule
\end{tabular}
\end{center}

\vspace{4pt}
\noindent Le passage de $1{,}380$ à $1{,}931$ pilier par sinistre, soit un facteur $1{,}40$, est ce
qui commande toute la suite : ce qui touche au coût des sinistres \emph{multi-piliers} frappe l'état
non conforme deux fois plus que l'état conforme, donc déplace l'\emph{écart} et pas seulement le
niveau.
% SOURCES-SCRIPTS: 74

\subsection*{Le relâchement de l'additivité, et la table complète des exposants}

La perte d'un sinistre touchant $k$ piliers vaut $\bigl(\sum_j X_j\bigr)k^{\,\theta-1}$. La
paramétrisation a deux propriétés qui la rendent lisible : $\theta=1$ redonne \emph{exactement} le
modèle publié, ce qui sert de contrôle et a été vérifié au centime sur les deux états ; et
$1^{\theta-1}=1$, donc elle laisse les sinistres mono-pilier là où ils sont et ne touche que la
simultanéité.

\begin{center}\footnotesize
\renewcommand{\arraystretch}{1.25}
\begin{tabular}{@{}c c c c c@{}}
\toprule
$\boldsymbol{\theta}$ & \textbf{Conforme} & \textbf{Non conforme} &
$\boldsymbol{\Delta_{\text{DORA}}}$ & \textbf{/ publié} \\
\midrule
$0{,}70$ & $5\,435$ & $15\,141$ & $9\,706$ & $0{,}687$ \\
$0{,}85$ & $5\,732$ & $17\,536$ & $11\,804$ & $0{,}835$ \\
$1{,}00$ & $6\,049$ & $20\,188$ & $14\,139$ & $1{,}000$ \\
$1{,}15$ & $6\,449$ & $23\,447$ & $16\,998$ & $1{,}202$ \\
$1{,}30$ & $6\,972$ & $27\,459$ & $20\,487$ & $1{,}449$ \\
\midrule
\emph{élasticité} & $0{,}39$ & $0{,}95$ & $1{,}20$ & \\
\bottomrule
\end{tabular}
\end{center}

\vspace{4pt}
\noindent La plage de $10\,781$~M\euro{} vaut $15$~fois le bruit de simulation de l'écart, mesuré à
$\pm 734$~M\euro{} sur quatre graines appariées. L'effet n'est donc pas un artefact de Monte-Carlo,
et c'est ce qui autorise à le qualifier de plus lourde des hypothèses structurelles non testées.
% SOURCES-SCRIPTS: 74

\subsection*{Pourquoi l'amplitude est posée et non estimée}

Deux mécanismes mettraient l'additivité en défaut, et ils tirent en sens \emph{contraire}. La
sous-additivité d'abord : une cellule de crise, une investigation forensique et une communication
servent plusieurs piliers à la fois, si bien que le total serait inférieur à la somme et que le
modèle publié \emph{sur}-estimerait. La super-additivité ensuite : plusieurs défaillances
simultanées saturent la capacité de remédiation, les délais s'allongent, les coûts unitaires
montent, et le modèle \emph{sous}-estimerait.

Aucun des deux n'est séparable sur les sources retenues, et la raison est celle exposée au
\S\ref{sec:ann-proxys} : le coût d'un sinistre \emph{daté} n'y est jamais observé mais reconstruit,
la chronologie de brèches portant les dates et la conversion en montants portant les coûts. Le coût
d'un sinistre \emph{multi-piliers} le serait moins encore, puisqu'il faudrait une ventilation par
domaine de contrôle qu'aucune source publique ne fournit.

C'est pourquoi l'amplitude $[0{,}70\,;1{,}30]$ est posée et non estimée, et pourquoi la
grandeur à retenir est l'\emph{élasticité} plutôt que la plage : une élasticité vaut pour toute
amplitude qu'un lecteur jugera plausible, une plage ne vaut que pour celle qu'on a choisie. C'est
aussi la seule ligne du registre des écarts du \S\ref{sec:pu-coherence} dont le \emph{sens} ne se
déclare pas, les sept autres ayant une direction connue.
% SOURCES-SCRIPTS: 74

% =====================================================================
\section{Deux corrections d'instrument : le tornado et l'objet de l'invariance}
\label{sec:ann-corrections-instrument}
% SOURCES-SCRIPTS: 76 77

Cette pièce documente deux corrections apportées à des résultats déjà publiés. Aucune ne renverse
une conclusion, toutes deux en changent la portée ou le statut.

\textbf{Le tornado comparait des longueurs de barre, pas des sensibilités.} Le tornado du
\S\ref{sec:sensibilite-propagation} balaie des plages d'inégale vraisemblance, et l'un de ses
leviers, le seuil, n'en est pas un : le déplacer réajuste la loi de queue, donc quatre paramètres à
la fois.

\begin{center}\footnotesize
\renewcommand{\arraystretch}{1.2}
\begin{tabular}{@{}l r r r@{}}
\toprule
\textbf{Composante} & \textbf{p80} & \textbf{p90} & \textbf{Effet signé} \\
\midrule
Le paquet, ce que publie le tornado & $12\,944$ & $4\,862$ & $-8\,081$ \\
\midrule
Seuil $u$ seul & $6\,628$ & $6\,948$ & $+320$ \\
Indice de queue $\xi$ seul & $18\,533$ & $4\,540$ & $\mathbf{-13\,993}$ \\
Échelle $\sigma$ seule & $4\,617$ & $9\,299$ & $+4\,682$ \\
Taux de dépassement $p_u$ seul & $7\,339$ & $4\,814$ & $-2\,525$ \\
\midrule
\emph{somme des quatre} & & & $-11\,517$ \\
\bottomrule
\end{tabular}
\end{center}

\vspace{4pt}
\noindent Le seuil pris seul ne déplace que $+320$~M\euro{} quand l'indice de queue seul en
déplace $-13\,993$ : la barre publiée sous le nom de \og seuil \fg{} mesure la sensibilité à la
queue. Les composantes se compensent au lieu de s'additionner, leur somme signée valant
$-11\,517$ pour un paquet à $-8\,081$. Lue en élasticités, à perturbation relative égale et un
paramètre à la fois, la conclusion tient : la propagation reste le plus faible des leviers
substantiels, $0{,}37$ contre $3{,}79$ pour l'indice de queue.

\textbf{L'invariance de la dépendance porte sur une espérance, pas sur un quantile.} L'argument
d'additivité du \S\ref{sec:dependance-etats} vaut pour une espérance, or le capital est un
quantile. Les deux objets ont donc été calculés : la moyenne des VaR conditionnelles aux $32$
configurations vaut $9\,739$~M\euro{} et le quantile de la loi mélangée $10\,643$, soit
$904$~M\euro{} ou $9{,}3\,\%$ d'écart, parce qu'à $99{,}5\,\%$ la queue du mélange est prélevée
dans les composantes les plus lourdes. Le mémoire publie la première grandeur, et le dire fait
partie du résultat. Sur le balayage, l'espérance ne bouge que de $2$~M\euro{} quand le quantile se
déplace de $334$, valeur non monotone et du même ordre que le bruit de tirage, dont l'étendue vaut
$344$~M\euro{} sur huit graines : l'invariance n'y est pas démontrée, elle est seulement non
détectée, ce qui est strictement plus faible. Ces $334$~M\euro{} valent $18\,\%$ de la largeur de
la bande d'identification, donc la conclusion de gestion ne change pas ; c'est son statut qui
change.

\begin{figure}[htbp]\centering
\figover{S33_var_melange.png}
\caption{Les deux objets. (a)~la moyenne des VaR conditionnelles et le quantile de la
loi mélangée diffèrent de $9{,}3\,\%$, ce qui impose de dire lequel est publié.
\textbf{(b)}~sur le balayage, l'espérance ne bouge pas et le déplacement du quantile est du même
ordre que le bruit de tirage, sans être monotone : il n'est pas résolu. L'échelle est celle de la
bande d'identification déjà publiée.}
\label{fig:var-melange}
\end{figure}

% =====================================================================
\section{Le plafond de sévérité, mesuré sur les quatre bilans réels}
\label{sec:ann-plafond}

% SOURCES-SCRIPTS: 78

Cette pièce accompagne la limite déclarée au \S\ref{sec:pu-coherence} sur l'absence de plafond de
sévérité à l'échelle d'entité. Elle donne la figure ; le raisonnement et les nombres sont dans le
chapitre.

\begin{figure}[htbp]\centering
\figover{S34_plafond_exposition.png}
\caption{Le plafond de sévérité adossé à l'exposition. (a)~la prédiction naturelle, le
quantile annuel borné par le plafond au titre du principe de la perte unique dominante, ne tient
qu'à plafond \emph{lâche} : plafonner détruit la queue lourde que ce principe suppose, et la
charge redevient un cumul. (b)~la première série donne la fraction du bilan qu'un sinistre
unique devrait détruire pour que le besoin publié soit atteignable, la seconde le plafond qui
ampute ce besoin de moitié. L'écart entre les deux barres est ce qui met la plus petite entité hors
domaine, et l'entité la plus grande n'est amputée par aucun plafond discutable.}
% SOURCES-SCRIPTS: 78
\label{fig:plafond-exposition}
\end{figure}

\paragraph{Une précaution de méthode, et elle vaut d'être lue.} Ce calcul simule au niveau du
\emph{sinistre}, alors que la chaîne d'entité du \S\ref{sec:ann-reconciliation} passe par
l'évaluateur d'identification partielle. Les deux machineries ne coïncident pas exactement, et
l'écart reste sous $7\,\%$ sur les cinq entités. Toutes les comparaisons de cette section sont
donc faites \emph{à l'intérieur d'un seul moteur}, chaque entité étant rapportée à sa propre
valeur non plafonnée, sans quoi l'effet du plafond se mélangerait à celui de la machinerie.

Une raison technique explique d'ailleurs que ce calcul n'ait pas été fait plus tôt : la chaîne
d'entité met la \emph{somme} des pertes à l'échelle par un multiplicateur de sévérité, ce qui est
exact pour une multiplication mais faux pour un plafond, plafonner chaque sinistre puis
sommer n'étant pas plafonner la somme. Il fallait redescendre au niveau du sinistre.

% =====================================================================
\section{Les deux horloges du modèle et le coût de leur séparation}
\label{sec:ann-horloges}
% SOURCES-SCRIPTS: 79

Le modèle sépare une horloge \emph{lente}, l'année qui porte la fréquence et sur laquelle le
capital se mesure, d'une horloge \emph{rapide}, la propagation à l'intérieur d'un sinistre. Deux
approximations en découlent, la cascade traitée comme instantanée et l'état de conformité figé sur
l'année, et cette section mesure ce que la commodité coûte.

\textbf{La cascade instantanée est une approximation neutre, et non prudente.} On donne une durée à
la cascade et l'on compare à nombres communs, deux valeurs de délai partageant les mêmes tirages à
état et graine fixés. Au délai que la donnée soutient, $6{,}7$~h par pas, l'écart de capital vaut
$14\,139$~M\euro{} contre $14\,139$ pour la cascade instantanée, soit un déplacement nul sur les
quatre graines : une cascade de quelques heures ne rencontre pratiquement jamais le 31 décembre. À
vingt-quatre heures par pas le déplacement reste sous le demi-pour-cent, et au-delà il change de
signe d'une graine à l'autre, deux mécanismes opposés que quatre graines ne séparent pas.

\textbf{Le capital est concave en la durée de non-conformité.} En faisant varier la part $\tau$ de
l'année passée en non-conformité, la courbe passe au-dessus de sa corde, de $897$~M\euro{} à
mi-année, soit $2{,}2$ fois le bruit de simulation. La signature est celle d'une queue portée par
un sinistre unique : il suffit d'une fenêtre de non-conformité pour qu'un tel sinistre s'y
produise.

\begin{cle}
\textbf{Deux lectures de gestion, et aucune n'est celle que l'intuition suggère.} Un trimestre de
non-conformité porte déjà $33\,\%$ du surcoût de l'année entière, donc une fenêtre courte n'est pas
une petite exposition ; et remédier à mi-exercice ne capture que $44\,\%$ du bénéfice annuel, non
la moitié, de sorte qu'un plan qui proratise le gain se trompe du côté optimiste. Le collapse de
l'horloge est donc légitime pour ce que le mémoire en fait, publier un capital par état, et fautif
dès qu'on veut en lire un calendrier de remédiation : la limite est d'usage et non de construction.
\end{cle}

\noindent Un lien manquait enfin avec l'additivité des coûts. Au délai mesuré, $8{,}0\,\%$ des
sinistres non conformes en recouvrent un autre contre $1{,}7\,\%$ à l'état conforme, soit un
facteur $4{,}8$, qui monte par deux canaux composés, la non-conformité multipliant le nombre de
sinistres et allongeant chaque cascade. Or l'exposant $\theta$ du \S\ref{sec:additivite-couts} ne
joue qu'au sein d'un sinistre : la plage publiée en sous-estime donc l'amplitude. L'énoncé porte
sur l'amplitude et non sur le sens, la concomitance amplifiant la mutualisation autant que la
saturation ; en revanche l'asymétrie entre états est renforcée dans le sens du mémoire, ce qui
montre que les élasticités de $0{,}95$ contre $0{,}39$ ne sont pas un artefact de $\theta$.

% =====================================================================
\section{Le plafond et la saturation, mesurés conjointement}
\label{sec:ann-plafond-saturation}
% SOURCES-SCRIPTS: 86

Deux réserves du mémoire déplacent le capital en sens contraire : le plafond de sévérité adossé à
l'exposition le fait baisser, la saturation des coûts relâchée par l'exposant $\theta$
(\S\ref{sec:additivite-couts}) le fait monter. Se compensent-elles, de sorte que l'une
dispenserait de discuter l'autre ? Le contrôle passe d'abord : sans plafond ni saturation, la VaR
de l'entité notionnelle vaut $179{,}1$~M\euro, du même ordre que les $177{,}1$ publiés au
\S\ref{sec:ann-plafond} et que les $169$ de l'évaluateur d'identification partielle. La saturation
seule fait passer la VaR de $144{,}2$ à $225{,}7$~M\euro{} sur la plage de $\theta$, soit une bande
de $81{,}5$ ; le plafond seul la fait tomber de $179{,}1$ à $56{,}9$~M\euro{} quand $\kappa$
descend à $0{,}2\,\%$, soit $122{,}2$.

\begin{cle}
\textbf{Le plafond est un amortisseur, non un contrepoids.} La bande conjointe vaut
$177{,}8$~M\euro{} contre $203{,}7$ pour la somme des deux bandes marginales, soit un rapport de
$0{,}87$ : les additionner surestimerait l'incertitude. Mais les deux réserves ne se neutralisent
pas. Le terme d'interaction change de signe à $\theta = 1$, et à $\kappa = 0{,}5\,\%$ la
sensibilité à $\theta$ est exactement nulle, le plafond absorbant alors entièrement la saturation.
Il n'existe aucune crête de compensation non triviale : les deux réserves sont hiérarchisées et non
confondues, et l'on peut discuter le plafond sans discuter $\theta$.
\end{cle}

% =====================================================================
\section{Les séquences ordonnées du corpus et leur absence de pouvoir identifiant}
\label{sec:ann-sequences}
% SOURCES-SCRIPTS: 75 59

Cette pièce accompagne le \S\ref{sec:sequences-ordonnees}. L'idée examinée est légitime, des marges
par paires ne déterminant pas une loi sur les permutations, mais elle ne donne rien ici, et le
motif vaut mieux qu'un constat de pénurie.

\textbf{Les séquences étaient déjà dans le corpus.} Chaque incident code un ensemble d'arêtes
dirigées, jamais une chaîne ; les composer produit $9$ chemins de longueur deux pour $4$ séquences
distinctes, la plus longue chaîne observée étant P1 $\to$ P5 $\to$ P3 $\to$ P2. Quatre incidents
sur dix n'en fournissent aucune.

\textbf{La comparaison demandée a un support vide, et c'est structurel.} Elle oppose
P1 $\to$ P2 $\to$ P3 à P2 $\to$ P3 $\to$ P1, et les deux exigent que P2 émette, or P2 n'émet
jamais, $0$ émission sur $22$ transitions contre $10$ réceptions ; symétriquement P1 n'est jamais
atteint. La comparaison n'est donc pas indécise faute d'échantillon, elle est sans objet. Le corpus
code par ailleurs le \emph{conditionnement direct} et non la \emph{précédence} : la clôture
transitive ajoute deux paires qu'aucun rapport ne code, donc lire $p_{jk}$ comme une précédence
serait une faute.

\textbf{L'acyclicité n'est pas un résultat.} Le graphe agrégé est acyclique, d'ordre topologique
$\{$P1$\} < \{$P4,\,P5$\} < \{$P3$\} < \{$P2$\}$, mais sous le nul par paire $42$ des $64$
orientations le sont aussi, soit $65{,}6\,\%$ : on obtiendrait l'acyclicité observée deux fois sur
trois par hasard. Le nul par observation la rendrait écrasante, mais il détruit la constance du
codage et mesure la fiabilité du codeur, non l'existence d'une direction.

\textbf{Le modèle prédit les séquences : il n'est pas sans mémoire.} La cascade étant
auto-évitante, la loi du successeur est renormalisée sur les piliers non visités, donc elle dépend
du prédécesseur. Le gonflement moyen vaut $1{,}211$ sur les quatre séquences observées, et c'est
une borne basse. Les triplets portent ainsi une prédiction testable, les chemins de longueur deux
pesant $18{,}29\,\%$ de la masse du modèle. Le rapprochement avec le corpus n'est toutefois pas un
test, celui-ci n'étant pas un échantillon tiré au hasard mais un ensemble de post-mortems
notables : tout ce qu'on peut en dire est qu'il n'y a pas de contradiction criante.

\textbf{Le test n'a aucune puissance, et la condition de réouverture est structurelle.} Il exige
d'un même pilier qu'il soit atteint et qu'il ait deux successeurs distincts. P1 a trois successeurs
mais n'est jamais atteint ; P3, P4 et P5 sont atteints mais leur loi de successeur est un point ;
P2 n'émet pas. Aucun pilier ne remplit les deux conditions et les deux manques sont exclusifs :
la statistique est identiquement nulle. Ajouter dix incidents du même genre laisserait les lois
dégénérées, donc le critère de réouverture porte sur la structure du corpus avant son volume. Le
volume nécessaire se chiffre ensuite sur les écarts que le modèle prédit, de $158$ à $589$ chemins
passant par un même pilier selon la séquence visée, pour séparer au seuil de $5\,\%$ avec $80\,\%$
de puissance ; le corpus en fournit $9$.

\begin{figure}[htbp]\centering
\figover{S31_sequences_ordonnees.png}
\caption{Ce que les séquences ordonnées disent, et ce qu'elles ne peuvent pas dire.
\textbf{(a)}~le graphe agrégé des $22$ transitions : épaisseur et chiffre donnent le nombre
d'incidents, P1 est source pure et P2 puits pur, et l'ordre est presque total.
\textbf{(b)}~les deux conditions d'un test de dépendance au prédécesseur, côte à côte. P1 a trois
cibles mais n'est jamais atteint ; les trois autres émetteurs sont atteints mais n'ont qu'une
cible. Les deux conditions ne se rencontrent sur aucun pilier.}
\label{fig:sequences-ordonnees}
\end{figure}

% =====================================================================
\section{La dynamique du facteur systémique et ses deux échelles d'aléa}
\label{sec:dynamique-Y}

Cette pièce accompagne la brique de fréquence du chapitre~\ref{chap:socle}. Le facteur commun
$Y_{j,t}$ y est posé centré et réduit ; sa \emph{dynamique} est écrite ici. La lacune n'est pas cosmétique : la loi de $Y_{j,t}$ ne change aucune probabilité
marginale de défaillance, mais elle gouverne entièrement la façon dont plusieurs piliers
tombent \emph{ensemble}, ce qui est l'objet même de ce mémoire.

\paragraph{Deux échelles, et un argument empirique qui interdit de n'en garder qu'une.}
Un système d'information subit en permanence un grand nombre de petites atteintes dont aucune
ne déplace seule un état de conformité : la somme normalisée de ces contributions converge, par
le théorème central limite sous la condition de Lindeberg, vers une gaussienne, et son
accumulation en temps continu est donc un mouvement brownien. Cette forme est imposée par la
structure de l'aléa. À l'autre extrémité se trouvent
les campagnes qui touchent simultanément de nombreuses entités par une vulnérabilité partagée,
et la condition de Lindeberg y échoue précisément parce qu'un seul événement domine la
somme de l'année : un brownien ne peut pas les représenter, ses accroissements étant sans
atome. Elles demandent une composante de sauts. La surdispersion mesurée, $1{,}19$ sur les
cellules firme-année et $2{,}24$ sur la série annuelle
% SOURCES-SCRIPTS: 47
détendancée, établit l'existence d'un facteur commun, mais elle ne départage pas les deux
% SOURCES-SCRIPTS: 08b
régimes : tout facteur commun non dégénéré, gaussien compris, produit un indice de dispersion
supérieur à un. Ce qui motive la composante de sauts est ailleurs, dans la dépendance de queue :
un facteur gaussien de corrélation inférieure à un donne un coefficient de dépendance de queue
nul, quand les campagnes de masse font céder de nombreuses entités ensemble. La seconde échelle
est donc une hypothèse de structure motivée par ces campagnes, non une conséquence de la
surdispersion.

\paragraph{L'écriture, et le paramètre interprétable qu'elle fait apparaître.}
Pour chaque pilier,
\begin{equation}
  \mathrm{d}Z_j(t) \;=\; -\kappa_j\,Z_j(t)\,\mathrm{d}t \;+\; \sigma_j\,\mathrm{d}B_j(t)
  \;+\; \mathrm{d}\!\Bigl(\textstyle\sum_{n\ge 1}\zeta^{(j)}_n\mathbf 1\{\tau_n\le t\}\Bigr)
  \;-\; \nu_j\mu_j\,\mathrm{d}t,
  \label{eq:sde-Y}
\end{equation}
où $(B_1,\dots,B_5)$ est un brownien à cinq composantes de corrélation instantanée portant la
co-occurrence \emph{ordinaire}, $(\tau_n)$ les dates d'un processus ponctuel d'intensité $\nu_j$
et d'amplitudes $\zeta^{(j)}_n$ pour les campagnes, le dernier terme étant le compensateur qui
garantit $\E[Z_j(t)]=0$. Le facteur qui entre dans la latente est l'état normalisé
$Y_{j,t}=Z_j(t)/\varsigma_j$ avec
\begin{equation}
  \varsigma_j^2 \;=\; \frac{\sigma_j^2}{2\kappa_j} \;+\;
  \frac{\nu_j\,\E\bigl[(\zeta^{(j)})^2\bigr]}{2\kappa_j},
  \qquad
  \pi_j \;=\; \frac{\nu_j\,\E\bigl[(\zeta^{(j)})^2\bigr]}
                   {\sigma_j^2+\nu_j\,\E\bigl[(\zeta^{(j)})^2\bigr]} \;\in\; [0,1],
  \label{eq:norm-Y}
\end{equation}
la variance stationnaire se lisant composante par composante, les deux sources étant
indépendantes. La part $\pi_j$ de variance systémique portée par les sauts est la grandeur
interprétable de cette écriture : nulle, le pilier ne subit que du bruit de fond ; égale à un,
il ne subit que des campagnes.

% SOURCES-SCRIPTS: 03
\begin{cle}
\textbf{Ce que cette écriture change, et ce qu'elle ne change pas.} Elle ne déplace
aucune probabilité marginale, donc aucun capital publié : le probit ne voit
$Y_{j,t}$ qu'à travers sa loi centrée réduite, et la normalisation \eqref{eq:norm-Y} impose
$\E[Y_{j,t}]=0$ et $\Var[Y_{j,t}]=1$ par construction, de sorte que l'incidence de fond reste
$\Phi(-1{,}7)$. Elle change la \emph{dépendance de queue}, et c'est tout son objet : le
coefficient de dépendance de queue supérieure d'une gaussienne de corrélation $r<1$ est
\emph{nul}, celui d'un mélange à choc commun est strictement positif. Et elle ne dégrade pas le
résultat d'identification partielle, la composante de sauts se manifestant par de la
co-occurrence, donc dans la partie symétrique de la matrice de propagation, et non dans sa
partie antisymétrique qui reste bornée.
\end{cle}

\paragraph{L'hypothèse retenue, déclarée plutôt que supposée.} On publie la forme à
renouvellement, $\kappa_j\to\infty$, où le facteur est l'accroissement normalisé de la période
sans mémoire d'un exercice à l'autre. Cela revient à poser $\theta_j=\kappa_j\times$ (durée
d'une grappe) $=\infty$, donc à supposer que le regroupement intra-annuel n'a aucun effet, et
cette hypothèse doit être écrite plutôt que laissée au lecteur. Trois motifs la soutiennent, et
ils ne se valent pas. Le premier est de fond : le regroupement est du second ordre sur le
quantile annuel par le principe du grand saut unique, la charge agrégée à queue régulièrement
variable étant portée par \emph{un} sinistre, et du premier ordre sur la co-occurrence entre
piliers, laquelle est déjà portée par la cascade et le choc commun du pilier des tiers. Le
mécanisme n'est donc pas ignoré, il est logé dans la coupe plutôt que dans le temps. Le
deuxième est d'identifiabilité, $\kappa_j$ n'étant pas estimable sur la profondeur d'historique
disponible. Le troisième est le gel de la calibration, et c'est le plus faible : il dit quand le
faire, pas s'il faut le faire.

\begin{attention}
\textbf{Le regroupement est réel, mesuré, et il ne mesure pas ce qu'on croit.} La charge
journalière du secteur financier a un indice de Gini de $0{,}69$ et une vague de mai 2023 porte
$191$ victimes en cinq jours ; l'excitation lue sur les dates au jour est
% SOURCES-SCRIPTS: 08g
intra-journalière, avec un pic à $3{,}3$~heures. Mais ces $191$ victimes
% SOURCES-SCRIPTS: 08h
sont $191$ \emph{entités distinctes} atteintes par une cause commune, non une entité subissant
$191$ sinistres : ce que le Gini mesure est de la dépendance en coupe, pas une
accumulation intra-annuelle pour une entité donnée. Or le besoin de capital calculé est celui
d'une entité, et à cette échelle la question est presque vide puisque la quasi-totalité des
années y sont sans aucun incident matériel. Le regroupement observé est un phénomène de
\emph{marché}, et c'est à ce titre qu'il entre dans le modèle.
\end{attention}

% =====================================================================
\section{Le quantile de la loi des configurations et la polarisation}
\label{sec:ann-quantile-configurations}

% SOURCES-SCRIPTS: 94

Cette pièce accompagne le \S\ref{sec:dependance-etats}.

\paragraph{Reste un troisième objet, et c'est celui que la réserve visait. Il est maintenant
calculé.} Les deux grandeurs qui précèdent portent sur la loi de \emph{perte} : la moyenne des
quantiles conditionnels, et le quantile du mélange. La réserve posée par l'argument d'additivité
visait autre chose, un quantile de la loi des configurations elle-même, c'est-à-dire de
l'incertitude sur l'\emph{état} de conformité de l'entité. Sur le même balayage, à marges
inchangées, l'espérance ne bouge que de $-0{,}02\,\%$, ce qui reproduit le résultat déjà publié,
tandis que le quantile de niveau $75\,\%$ de cette loi bouge de $+5{,}53\,\%$, celui de $90\,\%$
de $+2{,}76\,\%$ et la moyenne de queue au-delà de $90\,\%$ de $+1{,}06\,\%$. L'argument
d'additivité ne transporte donc pas à un quantile, et la réserve était justifiée. Mais l'ordre
de grandeur est celui de quelques pour cent, contre deux centièmes de pour cent pour l'espérance :
deux ordres de grandeur d'écart entre les deux objets, et les deux restent loin des incertitudes
déclarées ailleurs.

Le mécanisme est la polarisation, et c'est la mesure la plus robuste de cet exercice
parce qu'elle varie continûment quand un quantile discret, lui, procède par sauts. À espérance
inchangée, davantage de corrélation rend les configurations extrêmes plus probables des
\emph{deux} côtés : la probabilité que les cinq piliers soient conformes passe de $29{,}06$ à
$33{,}03\,\%$ et celle qu'ils soient tous non conformes de $6{,}83$ à $8{,}78\,\%$, tandis que la
probabilité de dépasser $1{,}25$ fois le capital espéré monte de $24{,}22$ à $27{,}34\,\%$.

\begin{attention}
\textbf{Deux mises en garde sur cet objet, et la première est une limite de l'objet lui-même.}
La loi des configurations ne porte que $32$ atomes, et le plus haut en concentre $6{,}83\,\%$ :
au-delà du niveau de $93{,}17\,\%$, son quantile \emph{est} la configuration intégralement non
conforme et ne peut plus bouger. Un \og{}quantile à $99{,}5\,\%$ de la loi des
configurations\fg{} n'aurait donc aucun contenu, et l'immobilité des niveaux hauts est une
\emph{saturation}, non une invariance. On ne demande pas à une loi à $32$ atomes la résolution
qu'on demande à une loi de perte.

\emph{Et cette grandeur n'est pas une mesure de capital.} Elle porte sur l'incertitude d'état,
non sur la charge annuelle : elle ne reçoit ni $\VaR$ ni $\TVaR$, réservées à cette dernière, et
la confondre avec l'exigence de solvabilité ferait lire une incertitude d'état comme une
exigence. C'est la même discipline de notation que celle qui sépare $\VaR$ de $\qsev$.
\end{attention}

% =====================================================================
\section{Le biais de narration : loi nulle exacte, jackknife et second codage}
\label{sec:ann-biais-narration}

% SOURCES-SCRIPTS: 64 54

Cette pièce accompagne le \S\ref{sec:biais-narration}.

\paragraph{Une loi nulle exacte.} L'asymétrie se
décompose sur les paires non ordonnées,
\[
  T \;=\; \tfrac12\lVert M - M^\top\rVert_1
    \;=\; \sum_{\{j,k\}} \bigl| m_{jk} - m_{kj} \bigr|,
\]
et sous $H_0$ chaque arête s'oriente indépendamment à pile ou face. Le solde d'une paire portant
$n$ arêtes suit donc la loi de $2B-n$ avec $B$ binomiale $(n,1/2)$, et les paires étant
indépendantes, la loi de $T$ s'obtient par convolution exacte, sans aucun tirage. On y
gagne deux choses. La $p$-valeur n'est plus plafonnée par le nombre de permutations : elle vaut
$1{,}53\cdot10^{-5}$, soit une chance sur $65\,536$, là où les $20\,000$ permutations butaient
sur $10^{-4}$. Et rejouer le test des milliers de fois ne coûte plus rien, ce qui est
exactement ce qu'exige la suite.

\begin{attention}
\textbf{Le $z$ cesse d'être le critère.} La loi nulle exacte est discrète et asymétrique. Dans
la configuration la plus défavorable examinée ci-dessous, on lit $z=+1{,}85$, donc au-dessus du
repère gaussien $1{,}645$, alors que la $p$-valeur exacte vaut $0{,}084$, donc au-dessus de
$0{,}05$. Le $z$ y \emph{surestime} la significativité. Toutes les décisions de cette section
sont donc prises sur la $p$-valeur exacte, le $z$ n'étant conservé que comme grandeur de
lecture comparable à celle des sections précédentes.
\end{attention}

\paragraph{Aucun rapport ne porte le signal à lui seul.} Le jackknife retire un incident à la
fois et rejoue le test exact. Le pire retrait, celui de TSB Bank 2018, laisse
$p = 2{,}44\cdot10^{-4}$ ; les dix retraits restent sous $3\cdot10^{-4}$. Ce n'était pas
surprenant vu la cohérence totale du corpus, mais ce n'est pas non plus ce qui le menace : le
danger n'est pas un rapport aberrant, c'est un biais \emph{systématique} présent dans tous les
rapports à la fois.

\paragraph{Combien le biais devrait-il fabriquer pour renverser la conclusion ?} On suppose donc
que la convention narrative a produit une partie des arêtes partant de $P_1$, et l'on cherche
l'ampleur qu'il faudrait. Deux opérations, la seconde bien plus hostile : le retrait,
où le rapport n'établissait rien, et le retournement, où l'arête allait en réalité dans
l'autre sens. Dans les deux cas on énumère les $2^{11}$ sous-ensembles d'arêtes et on retient le
\emph{pire}, jamais un tirage moyen : c'est une borne adverse, pas une espérance.

\begin{attention}
\textbf{La quatrième ligne ne passe pas, et il faut la donner quand même.} Retirer les sept
\emph{rapports} où $P_1$ apparaît comme source, au lieu des seules arêtes, ne laisse que trois
transitions et $p = 0{,}50$. Cette variante jette cependant avec l'eau du bain des arêtes que le
biais ne concerne pas, les chaînes $P_3\to P_2$ documentées dans les \emph{mêmes} rapports : un
biais sur la désignation de la cause racine n'invalide pas la séquence technique que le même
rapport établit par ailleurs. La borne pertinente est celle du retrait des arêtes ; celle du
retrait des rapports est une borne de pure forme, et elle est publiée ici pour que personne
n'ait à la reconstituer contre le travail.
\end{attention}

Le dispositif est préparé et disponible, et c'est à ce titre qu'il est décrit ici. Le
kit couvre les dix récits du corpus étendu, chacun réduit aux faits établis par le
rapport, sans aucune des flèches du modèle, et les entités y sont anonymisées pour que le codeur ne
reconnaisse pas le cas et ne se souvienne pas de la cause racine publiée. La grille qu'il remplit
porte les mêmes étiquettes anonymes que les récits : la nommer aurait rompu l'aveuglement par le
fichier de réponse lui-même. Sur les $10 \times 10 = 100$ cellules, le codage de référence en
compte $22$ directionnelles, ce qui fixe la base d'accord sur laquelle le $\kappa$ de Cohen se
calcule. En l'absence de second jeu, le script ne fabrique aucune donnée et s'arrête :
un dispositif qui produirait un coefficient sans second codeur vaudrait moins que son absence.

\begin{figure}[htbp]
  \centering
  \includegraphics[width=0.99\textwidth]{Z22_biais_narration.png}
  \caption{Le biais de narration, mesuré au lieu d'être déclaré.
\textbf{(a)}~jackknife par incident : aucun rapport ne porte le signal à lui seul, les dix
retraits laissant $p \leq 2{,}4\cdot10^{-4}$. (b)~point de rupture, en $p$-valeur
exacte et non en $z$ : il faudrait quatre arêtes de $P_1$ retournées dans la pire configuration
pour perdre la significativité, et aucun nombre de retraits n'y suffit. La courbe de
retournement remonte au-delà de sept parce que retourner \emph{toutes} les arêtes n'efface pas
l'asymétrie, elle l'inverse. (c)~le niveau de capital résiste, y compris avec $P_1$
rendue muette, et la bande d'identification se resserre à mesure que l'émission de $P_1$
disparaît.}
  \label{fig:biais-narration}
\end{figure}
% SOURCES-SCRIPTS: 64

% =====================================================================
\section{La direction bayésienne : prior sceptique et crédibilité par paire}
\label{sec:ann-bayes}

% SOURCES-SCRIPTS: 56

Cette pièce accompagne le \S\ref{sec:bayes}.

Deux vérifications s'imposent, et elles vont dans des sens opposés. La première est rassurante :
sous un prior sceptique $\mathrm{Beta}(4,4)$, qui tire vers $\theta=1/2$ donc vers
l'absence de direction et qui est le plus défavorable à la thèse, la médiane ne se déplace que de
$0{,}8\,\%$. La conclusion vient donc des comptes documentés, non du prior. La seconde est une
mise en garde, et elle est décisive : par paire, seules deux des six paires documentées
ont un intervalle de crédibilité à $90\,\%$ qui exclut $\theta=1/2$. Avec une à trois chaînes
observées, l'espérance penche nettement mais la crédibilité ne suffit pas. Le test de signe de
la section~\ref{sec:postmortems} répondait à une \emph{autre} question, celle de la cohérence
d'ensemble des dix-sept transitions ; le posterieur juge chaque paire séparément et se montre
plus prudent. Les deux lectures sont justes, et c'est la plus conservatrice qui doit être
retenue.

\begin{figure}[htbp]
\centering
\figover{Z19_bayes_direction.png}
\caption{Reformulation bayésienne de la direction. (a)~posterieurs par paire : les paires
documentées s'écartent de $\theta=1/2$, celles sans donnée restent uniformes. (b)~le capital
devient une loi, plus informative que des bornes de pire cas. (c)~robustesse au prior : la
médiane bouge de $0{,}8\,\%$, même sous prior sceptique.}
% SOURCES-SCRIPTS: 56
\label{fig:bayes}
\end{figure}

% =====================================================================
\section{Le benchmark copule sur plusieurs graines et plusieurs niveaux}
\label{sec:ann-benchmark-graines}

Cette pièce accompagne le \S\ref{sec:benchmark}.

\textbf{Le $+0{,}1\,\%$ est une graine, et la coïncidence est plus fragile que le nombre ne le
laisse croire.} La reprise du protocole sur quatre graines et un échantillon assez grand pour
qu'un quantile à $99{,}95\,\%$ ait un sens donne, au même niveau
% SOURCES-SCRIPTS: 80
$99{,}5\,\%$, un écart moyen de $+4{,}8\,\%$ d'étendue $10$ points entre graines. Le
$+0{,}1\,\%$ publié n'est donc pas faux, il est \emph{un tirage} : la conclusion à retenir n'est
pas que la gaussienne retombe exactement sur la cascade, c'est que leur écart ne dépasse
pas son propre bruit. L'énoncé faible est le seul qui tienne, et il suffit à l'argument.

\textbf{Et la séparation en queue ne vaut que pour la Student.} Le même protocole, décliné sur
huit niveaux de quantile de $50$ à $99{,}95\,\%$, marque d'une étoile les seules cases où l'écart
dépasse l'étendue entre graines. Au-delà de $99\,\%$, seule la Student sort de son
bruit, à $+20{,}4$, $+29{,}8$ puis $+31{,}2\,\%$ ; ni la gaussienne ni l'indépendance n'y
parviennent, leurs étendues atteignant $13$ et $31$ points. La raison est structurelle : la
Student est la seule des trois à porter une dépendance de queue \emph{asymptotique}. Le préprint
doit donc être cité sur le protocole, jamais sur l'ordre du résultat en queue, qui
dépend de l'indice de queue de la sévérité et non de l'architecture du modèle.

% =====================================================================
\section{La précision de simulation des six postures}
\label{sec:ann-postures}

% SOURCES-SCRIPTS: 46 51

Cette pièce accompagne le \S\ref{sec:predictive-robuste}.

\textbf{Les deux extrémités de l'échelle sont exactes, tout le bruit est au milieu.} Le point est
le quantile d'un ajustement unique et le pire-cas celui d'un $\xi$ \emph{posé} : ni l'un ni
l'autre ne comporte de simulation. Ce sont les postures intermédiaires, celles qui intègrent
l'incertitude, qui en héritent un bruit propre. Une posture plus riche est donc, à ce titre
précis, moins reproductible.

\textbf{La posture la plus prudente est la moins précise.} À $\beta=99\,\%$ l'écart-type de
simulation vaut $24$~M\euro, environ le double de celui des deux niveaux inférieurs, et la raison
est structurelle : un percentile à $99\,\%$ de la loi bootstrap s'estime sur bien moins de
tirages utiles que sa médiane, dans une queue lourde. Cela ne disqualifie pas la posture, mais
cela interdit de la citer à trois chiffres significatifs. Sur quatre rééchantillonnages, un
écart-type est lui-même connu à $40\,\%$ près : le classement de $\beta=90\,\%$ et
$\beta=95\,\%$ n'est donc pas résolu ici et n'a pas à l'être, seul le saut vers $\beta=99\,\%$
est net.

\textbf{Une conséquence de rédaction.} Toute grandeur simulée de cette section est publiée avec son
bruit. La VaR prédictive vaut ainsi $648$~M\euro{} sur un ensemble bootstrap de $2\,000$ tirages et
$645$ sur un ensemble de $3\,000$, celui de la figure : même estimateur, et un écart de
$3$~M\euro{} inférieur à un écart-type de simulation. Les deux valeurs sont un seul nombre, et le
bruit publié dit la précision à laquelle le lire.

\begin{figure}[htbp]\centering
\figover{J2_var_predictive.png}
\caption{VaR prédictive. (a)~le paramètre $\xi$ est incertain. (b)~les queues
\emph{plug-in} et prédictive ne divergent que loin dans la queue. (c)~le capital se lit
en profil et en bande.}
% SOURCES-SCRIPTS: 46
\label{fig:predictive}
\end{figure}
